
% Authoritative LaTeX manuscript; Markdown is a generated compatibility projection.

% <!-- part-title: DocAble Among Other Software Factories -->

% <!-- chapter-title: Coda: Engineering the Factory -->

% <!-- coda: true -->

\chapter{Coda: Engineering the Factory}

% <!-- point: process-design-and-tolerance-describe-the-differences-without-a-maturity-model | The cases expose different arrangements of knowledge, production, evaluation, and authority rather than stages of a common maturity sequence. | terms: tolerance -->

The factories in this chapter arrange engineering knowledge, production, evaluation, and authority in different ways. DocAble provides a close account of how those arrangements developed within one project; the public cases show other choices under different organizational conditions. Some preserve obligations in explicit representations, some organize access to existing knowledge, and some change how work reaches review and admission. Their reported outcomes remain uneven, but their mechanisms make the choices concrete. They give engineers alternatives to examine, rather than stages of a common maturity sequence.

% <!-- point: supervision-capacity-and-carrying-cost-become-visible-as-realization-gets-cheaper | Process design and tolerance are complementary responses to differing organizational conditions; finite supervision and the upkeep of representations influence where engineers place control. | terms: tolerance -->

Process design and tolerance help explain those choices (\autoref{process-design-and-tolerance}). DocAble\textquotesingle s accessibility obligations and one-engineer staffing made explicit requirements and automated rejection especially valuable: I had little inherited human process to accommodate and no spare reviewers for every implementation decision. Coordinating agents still required work orders, scheduling, and completion checks. In an established organization, the factory must also fit existing teams, knowledge, and authority. Both settings require judgments about which properties to preserve and how to organize work around them. A more capable agent may change the process needed to produce a system without changing its memory bound or deployment restriction. Engineers can leave implementation choices open within those bounds, while weighing the cost of supervision against the cost of maintaining the representations and checks that support it.

% <!-- point: process-design-and-tolerance-give-machinery-for-the-theory-chapter | Human incentives and accountability remain distinct from agent capability; engineering authority determines delegation and controls, and the theory chapter examines their costs and effects. | terms: tolerance -->

Capability also does not settle accountability. Organizations delegate consequential judgment to people partly through incentives and responsibility, including professional obligations and institutional consequences\citep{marchsimon1958}. An agent does not bear those responsibilities or respond to those incentives in the same way a human engineer does. Its instructions and permissions can be changed, and its work can be stopped, but those actions do not transfer human accountability to it. Engineering authority must therefore decide which judgments to delegate, what evidence to require, and which obligations need controls. \magevalue{chapter:theory} examines the costs and effects of these choices: how representation and independent evaluation change the economics of realization, and how experience changes the engineering structures worth maintaining.
